\subsection{The Factor Theorem} 
This theorem \emph{seems} like something we already knew, but it actually says more.
\begin{theorem}[Factor Theorem]
For any polynomial, a number $r$ is a root of the polynomial if and only if $(x-r)$ is a factor of the polynomial.
\end{theorem}
We already knew that we can use the factors to find the roots, using our method of solving a polynomial equation by factoring. This theorem tells us that we can also use the roots to find factors. The phrase ``if and only if'' means that this theorem goes in both directions.
\begin{proof}
Call our polynomial $P(x)$. First, we'll show that if $(x-r)$ is a factor, then $r$ is a root.
\begin{itemize}
\item If $(x-r)$ is a factor, that means that there's some other polynomial such that
\[
\makeblank[3in]
\]
\item But then $P(r)=\makeblank[1.5in]=\makeblank[1.5in]=\makeblank[.5in]$.
\end{itemize}
Now we'll show that if $r$ is a root, $(x-r)$ is a factor.
\begin{itemize}
\item If $r$ is a root of the polynomial, $P(r)=\makeblank[.5in]$.
\item According to the \makeblank[1.5in] theorem, $P(x)=Q(x)(x-r)+R$, where $R=\makeblank[.5in]$
\item So $P(x)=Q(x)(x-r)+\makeblank[.5in]=\makeblank$
\item \makeblank[5in]\qedhere
\end{itemize}
\end{proof}
The point of this theorem is:
\begin{itemize}
\item If we know the \makeblank[1.5in] we can find the \makeblank[1.5in]
\item If we know the \makeblank[1.5in] we can find the \makeblank[1.5in]
\end{itemize}
We can use this to write down a polynomial given its roots.